%%%%%%%%%%%%%%%%%%%%%%%%%%%%%%%%%%%%%%%%%
% MODELE DE MANUSCRIT DE THESE 
% DE L'UNIVERSITE PARIS CITE
% OFFICIEL

% LaTeX Template
% Version overleaf (14/10/2023) 
% Format BOOK
% Par Aleksandra SAVINA, PhD
%%%%%%%%%%%%%%%%%%%%%%%%%%%%%%%%%%%%%%%%%


%----------------------------------------------------------------------------------------
%	PREAMBULE // PACKAGES AND OTHER DOCUMENT CONFIGURATIONS
%----------------------------------------------------------------------------------------
\documentclass[a4paper,12pt,twoside,french]{book}
[a4paper,12pt,twoside]
\usepackage[utf8]{inputenc}
\usepackage[T1]{fontenc}
\usepackage{babel}
%	MARGIN SETTINGS
\usepackage[a4paper,left=3cm,right=2.5cm,top=3cm,bottom=3cm, twoside]{geometry}

\usepackage{xcolor}
\usepackage{stmaryrd}
\usepackage{amssymb}
\usepackage{amsmath, amsthm}
% 创建一个定理样式
\newtheorem{thm}{Theorem}[section]
\newtheorem{defn}[thm]{Definition} % 共享编号

\usepackage{libertine}

\usepackage[most]{tcolorbox}

%\usepackage{cite}
\usepackage{hyperref} % for references (\ref, \label), url
\hypersetup{
    colorlinks=true,
    linkcolor=black,
    filecolor=black,      
    urlcolor=black,
    citecolor=black,
    pdfpagemode=FullScreen,
    }
\usepackage{cleveref}
\usepackage[
    backend=biber, 
    natbib=true,
    style=numeric,
    sorting=none
]{biblatex} %Imports biblatex package

\addbibresource{biblio.bib} %Import the bibliography file

\usepackage{graphicx} % to include images
%\usepackage[pdftex]{graphicx}
\graphicspath{ {figures/} }
\usepackage{caption}
\usepackage{subcaption}
\usepackage{float}
\usepackage{multirow}
\usepackage{array, multirow, tabularx}
\usepackage{booktabs}
\usepackage{longtable}
\usepackage{apalike}
\usepackage{minitoc}
\usepackage{tikz}
\usepackage{fancyhdr}
\pagestyle{fancy}
\usepackage[strict]{changepage}
\usepackage{siunitx}
\selectlanguage{french}%

\usepackage{enumitem}
\usepackage{epigraph}
%\usetikzlibrary{shadows.blur}
\usepackage{lscape}
\usepackage{titletoc}
\usepackage{mdframed}

\usepackage{tabu}
\usepackage{pdfpages}

\usepackage{arydshln}
%\usepackage[nonumberlist]{glossaries}
\usepackage{calc}
\usepackage[]{titlesec} 
\definecolor{linkColor}{HTML}{32a852}
%\usepackage[colorlinks=true,citecolor=linkColor,linkcolor=black]{hyperref}

\addto\captionsfrench{%
  \renewcommand{\listfigurename}{Liste des figures}%
}


%%%%%%%%%%%%%%%%%%%%%%%%%%%%%%%%%%%%%%%%%%%%%%%%%%%%%%%%%%%%%%%%%%%%%%%%%%
%%%%%%%%%%%%%%%%%%%%%%%%%%%%%%%%%%%%%%%%%%%%%%%%%%%%%%%%%%%%%%%%%%%%%%%%%%

\usepackage{sectsty}
%\sectionfont{\color{CentraleRed}}  % sets colour of sections
%\subsectionfont{\color{uclablue}}  % sets colour of subsections, ou gray, teal 
\definecolor{darkseagreen}{rgb}{0.56, 0.74, 0.56}
\usepackage{lipsum}

\title{NMR global data Version2}
\author{Aleksandra Savina}
\date{October 2022}

\newtcbtheorem[number within=chapter]{Theorem}{}{
        enhanced,
        sharp corners,
        attach boxed title to top left={
            xshift=-1mm,
            yshift=-5mm,
            yshifttext=-1mm
        },
        top=1.5em,
        colback=white,
        colframe=blue!75!black,
        fonttitle=\bfseries,
        boxed title style={
            sharp corners,
            size=small,
            colback=blue!75!black,
            colframe=blue!75!black,
        } 
    }{thm}

\newtcbtheorem[number within=chapter]{Definition}{}{
        enhanced,
        sharp corners,
        attach boxed title to top left={
            xshift=-1mm,
            yshift=-5mm,
            yshifttext=-1mm
        },
        top=1.5em,
        colback=white,
        colframe=blue!75!black,
        fonttitle=\bfseries,
        boxed title style={
            sharp corners,
            size=small,
            colback=blue!75!black,
            colframe=blue!75!black,
        } 
    }{def}

\newenvironment{myTheorem}[2]{ \begin{Theorem}[adjusted title=#1]{}{#2} 
  \textbf{Theorem \thetcbcounter.} }{\end{Theorem}}

  \newenvironment{myDefinition}[2]{ \begin{Definition}[adjusted title=#1]{}{#2} 
  \textbf{Definition \thetcbcounter.} }{\end{Definition}}

%----------------------------------------------------------------------------------------
%	HEADERS AND FOOTERS
%----------------------------------------------------------------------------------------

%\fancyhead[RO]{\scshape \leftmark}
%\fancyhead[L]{\scshape \righttmark}
\fancyhead{} % clear all header fields
\fancyfoot{} % clear all footer fields
%\fancyhead[LE,RO]{\slshape \rightmark}
\fancyfoot[C]{--~\thepage~--}

\usepackage{emptypage}

%\leftmark : adds name and number of the current top-level structure (for example, Chapter for reports and books classes; Section for articles ) in uppercase letters.
% \rightmark : adds name and number of the current next to top-level structure (Section for reports and books; Subsection for articles) in uppercase letters.

%%%%%%Glossaire
%\usepackage{glossaries}
\usepackage[acronym,xindy,toc]{glossaries} % nomain, if you define glossaries in a file, and you use \include{INP-00-glossary}
%[nomain,acronym,xindy,toc]
\makeglossaries

\usepackage[xindy]{imakeidx}
\makeindex
%\input{auxilliaires/glossaire}
%%%%%%%%%%%%%%%%%%%%%%%%%%%%%%%%%%%%%%%%%%%%%%%%%%%%%%%%%%%%%%%%%%%%%%%%%%
%%%%%%%%%%%%%%%%%%%%%%%%%%%%%%%%%%%%%%%%%%%%%%%%%%%%%%%%%%%%%%%%%%%%%%%%%%
%----------------------------------------------------------------------------------------
%	DOCUMENT
%----------------------------------------------------------------------------------------
\begin{document}

%\maketitle
%----------------------------------------------------------------------------------------
%	TITLE PAGE
%----------------------------------------------------------------------------------------
	\begin{titlepage}
		\begin{center}
			\begin{tabular}{c@{\hskip 7cm}c@{\hskip 1cm}}
				\includegraphics[height=2.5cm]{figures/PageTitre/UParisCitelogo.jpg} &
				%si autre logo : \includegraphics[height=2.3cm]{figures/PageTitre/logo_ic_2021.jpg}
			\end{tabular}
		\end{center}
	
		\begin{center}
		
\vspace*{.03\textheight}
\textsc{\LARGE Université Paris Cité}\\[0.2cm] % Univ name
		
        %Équipe de recherche\\
  
  			\vfill
 
	 		\rule{\textwidth}{0.8pt} \\ % Horizontal line
	 		\vspace{10pt}
	 		 { \LARGE \bfseries Interest rate and term structure modeling \textit{}} % Thesis title
	 		 \vspace{10pt}
	 		 \rule{\textwidth}{0.8pt} \\ % Horizontal line
		\end{center}
		
		\vfill
		% Author and supervisor	
		\begin{center}
			Presented by \textsc{\Large Hengze WANG / Ruoqi ZHU}\\[1cm] 
			Master project in \textsc{\large \large ISIFAR}\\[1cm]
		
	
			
			\vspace{1cm}
	
			April 2024
		\end{center}
		
		\vspace{1cm}

	\newpage %page de garde vide
	\thispagestyle{empty} %page de garde vide
	\end{titlepage}

\renewcommand{\chaptermark}[1]{\markboth{\textsc{#1}}{}}
\renewcommand{\chaptermark}[1]{\markboth{#1}{}}


\frontmatter % numérote les pages en chiffres romains jusqu'a la commande mainmatter
%\setcounter{secnumdepth}{3} % organisational level that receives a numbers
%\setcounter{tocdepth}{3}
\tableofcontents % Prints the main table of contents
\clearpage
%	Liste des abréviations-
%\chapter*{Liste des abréviations}
%\addcontentsline{toc}{chapter}{Liste des abréviations} 

\printglossary[type=\acronymtype, title=Liste des abréviations,toctitle=Liste des abréviations]

\newglossaryentry{latex}
{
        name=latex,
        description={Is a mark up language specially suited for 
scientific documents}
}

\begin{table}[h]
\centering
\caption{Table des Définitions des Symboles}
\begin{tabular}{|>{\bfseries}l|p{10cm}|}
\hline
Symbole & Définition \\
\hline
\( \Delta \) & Longueur de chaque sous-intervalle dans l'intervalle de temps \([0, \overline{T}]\) \\
\( \overline{T} \) & Longueur totale de l'intervalle de temps considéré \\
\( \overline{N} \) & Nombre de sous-intervalles dans lesquels l'intervalle de temps est divisé \\
\( t_n \) & Dates spécifiques où les transactions ont lieu, définies par \( t_n = n\Delta \) \\
\( p(n, N) \) & Prix d'une obligation T au temps \(t_n\) avec une maturité \(T = N\Delta\) \\
\( T \) & Temps de maturité de l'obligation T, qui est également \(N\Delta\) \\
\( p(N, N) \) & Prix de l'obligation T à maturité, par définition égal à 1 \\
\( p(0, N) \) & Famille de prix des obligations T aujourd'hui, varie avec \(N = 1, 2, \ldots, \overline{N}\) \\
\( p^*(0, N) \) & Structure à terme observée sur le marché pour les obligations à coupon zéro, représentant les données initiales pour les processus de prix \\
\([n-1, n]\) & Notation pour la période \(n\), souvent référencée comme l'intervalle de \(t_{n-1}\) à \(t_n\) \\
\hline
\end{tabular}
\end{table}




% --------------------------------------------------------------
%                         Début du corps
% --------------------------------------------------------------

\mainmatter
\setlength{\parskip}{.7em}

%\titlespacing*{\section}{0pt}{.9em}{.8em}
\renewcommand{\baselinestretch}{1.1}

\fancyhead[RO]{\leftmark}
\fancyhead[LE]{\textsc{\chaptername~\thechapter}}


\fancyhead{} % clear all header fields
\fancyhead[OL]{\textsc{Introduction}}
\chapter*{Introduction} % Main chapter title
\addcontentsline{toc}{chapter}{Introduction}  

À l'aide du livre «Financial mathematics : theory and problems for multi-period models»\cite{pascucci_financial_2012}, on va étudier le taux d'intérêt et ses modèles.

Les taux d'intérêt et la structure des échéances jouent un rôle central dans le fonctionnement des marchés financiers et dans les décisions de politique monétaire. Comprendre leur comportement et les modéliser de manière précise est crucial pour les investisseurs, les gestionnaires de risques, les décideurs politiques et les praticiens du domaine financier. Ce rapport se propose d'explorer en profondeur ces concepts fondamentaux, en mettant l'accent sur les méthodes de modélisation qui permettent de saisir la complexité et la dynamique des marchés des taux d'intérêt.

La modélisation des taux d'intérêt est une discipline interdisciplinaire qui fait appel à la théorie financière, à l'économétrie, aux mathématiques financières et à la statistique. Elle vise à expliquer et à prédire les variations des taux d'intérêt à différentes échéances, ainsi qu'à comprendre les relations entre ces taux et d'autres variables économiques et financières.

Au cours des dernières décennies, un large éventail de modèles a été développé pour capturer les caractéristiques complexes des taux d'intérêt et des structures d'échéances. Ces modèles vont des approches traditionnelles telles que le modèle de taux d'intérêt à terme à des modèles plus sophistiqués basés sur la théorie des probabilités, tels que les modèles à facteurs ou les modèles à sauts.

Dans ce rapport, nous commencerons par la définition de marché des taux d'intérêt. Nous présenterons les taux d'intérêt en introduisant les T-bonds et nous classerons les taux d'intérêt en différentes catégories sur la base de différentes propriétés et caractéristiques.

Ensuite, nous détaillerons deux types de modèles stochastiques d'évolution des taux d'intérêt, à savoir les modèles de type "short" et "forward". Pour le modèle short, nous intéressons au modèle affine et au modèle Hull-White. Pour le modèle forward, nous intéressons au modèle binomial forward. La consruction des ces modèles est pour le but d'identifier les différentes grandeurs économiques associées aux taux d'intérêt. %partie introduction ou avant-propos

\fancyhead{} % clear all header fields

\fancyhead[RO]{\leftmark}
\fancyhead[LE]{\textsc{\chaptername~\thechapter}}

\chapter{Marché de taux et les différents types de taux} % Main chapter title

%----------------------------------------------------------------------------------------
%	SECTION 
%----------------------------------------------------------------------------------------

Nous commençons par les principales caractéristiques des taux. Dans la suite, \( t < T < S \) désignent trois instants, un taux d'intérêt \( r \), relatif à l'intervalle de temps \([T,S]\), peut être de l'un des types suivants :

\textbf{Simple ou composé}:

\( r \) est simple si : 
\begin{equation}
I = C \times nr
\end{equation}

Où : 
\( I \) est les intérêts,
\( C \) est le capital initial qui s'exprime la somme d'argent initialement prêtée ou investie,
\( n \) est la durée pour laquelle l'argent est prêté ou investi, souvent exprimée en années.

\( r \) est composé si : 
\begin{equation}
A = P \times (1 + r)^n
\end{equation}

Où :
\( A \) est le montant accumulé après \( n \) périodes,
\( P \) est le principal ou le montant initial investi,
\( r \) est le taux d'intérêt par période,
\( n \) est le nombre de périodes d'investissement.

\textbf{Annualisé}:
\( r \) est un taux annualisé s'il est évalué sur une base annuelle.

\( r \) est un taux composé annualisé si : 
\begin{equation}
C_S = C_T e^{r(S-T)}
\end{equation}

où \( C_t \) représente la valeur du capital au moment \( t \)

\( r \) est un taux simple annualisé si : 
\begin{equation}
C_S = C_T(1 + (S - T)r)
\end{equation}

\textbf{Spot ou forward}:
\( r \) est un taux spot s'il est évalué à \( T \), c'est-à-dire au début de l'intervalle de référence.
\( r \) est un taux forward s'il est évalué à un instant \( t < T \), c'est à dire avant l'intervalle de référence.

\textbf{Taux forward annualisé simple} \( L(n;N,M) \), pour la période \([t_N, t_M]\), évalué en \( n \), d'après :
\begin{equation}
L(n;N,M) = \frac{1}{(M-N)\triangle} \left(\frac{p(n,N)}{p(n,M)} - 1\right)
\end{equation}

où :
    $\triangle=\frac{\bar T}{\bar N}$, et $t_n=n\triangle$ avec $0\le n \le \bar N$.
    
    $p(n,N)$ est le prix de $T$-bonds à l'instant $t_n=n\triangle$, de maturité $T=N\triangle$.
    
    Où $$D(n,N) := - \sum_{k = n}^{N-1} \exp{r^k} ,\quad 0 \le n \le N \le \bar{N}$$

Appelé facteur d'escompte sur la période $[t_n,t_N]$ : tous deux représentent la valeur au temps n d'une unité monétaire livrée au temps $N$. Cependant $p(n,N)$ est une valeur observable au temps $n$ car elle représente le prix d'un contrat qui est négocié sur le marché au temps $n$ ; par contre $D(n,N)$ n'est pas connue au temps $n$, car il s'agit d'une valeur aléatoire qui dépend de l'évolution des taux jusqu'à l'échéance. Cette observation sera rendue rigoureuse en termes mathématiques dans la section suivante.

\textbf{Taux forward composé} \( R(n;N,M) \), pour la période \([t_N, t_M]\), évalué en \( n \), d'après :
\begin{equation}
R(n,N) = R(n;N,N+1) = \log \left(\frac{p(n,N)}{p(n,N+1)}\right)
\end{equation}

$r_n := R(n,n)$ \textbf{taux spot composé}, pour la période $[t_n, t_{n+1}]$ : Nous appellerons simplement $r$ le taux court. Notons que, par définition, nous avons

\begin{equation}
p(n,n+1)= e^{-r_n}
\end{equation}

Notons $B$ la valeur du compte du marché monétaire, $B$ est donnée par la formule récursive : 
\begin{equation}
B_{n+1} = B_{n}e^{r_n}
\end{equation}

ou plus gégéralement : 
\begin{equation}
B_{N} = B_{n}\exp{\left(\sum_{k = n}^{N-1}r_k\right)}, \quad 0 \le n \le N
\end{equation}

où, par convention, nous supposons que $B_0 = 1$. Donc on a : 
\begin{equation}
B_{N} = \exp{\left(\sum_{k = 0}^{N-1}r_k\right)}
\end{equation}

\chapter{Modèles de marché pour les taux d’intérêt} % Main chapter title

\( r_n \) désigne une valeur aléatoire qui devient connue au temps \( n \), en particulier, le taux d'intérêt était supposé déterministe  même s'il n'était pas constant, à partir du paragraphe suivant, \( r \) (et par conséquent \( B \)) sera décrit par un processus stochastique. 

Il existe plusieurs approches de la modélisation stochastique des taux d'intérêt en temps discret. Nous examinons ici deux des principales classes de modèles, appelées respectivement de type court et de type forward. 
\section{Modèles Short}
Dans un modèle court, la dynamique du taux court \( r \) est définie par un processus stochastique approprié. Les obligations T, notées \( p(\cdot, N) \), sont considérées comme des dérivés du sous-jacent \( r \) avec une échéance \( N \) et un paiement final égal à 1, sur la base de la condition \( p(N, N) = 1 \). L'idée est d'utiliser la théorie classique de l'arbitrage pour obtenir les prix des obligations T à partir de la dynamique de \( r \). Voici quelques aspects de cette approche :

\begin{enumerate}
    \item Le taux court n'est pas un actif qui est négocié sur le marché : par conséquent, un modèle du type court est généralement, même dans les cas les plus simples, un modèle de marché incomplet, à savoir que la mesure de martingale n'est pas unique;
    \item Étant donné que la structure à terme initiale \( p^*(0, N) \), pour \( N = 1, \ldots, \bar{N} \), comprend des données de marché qui doivent être reproduites par le modèle, il est nécessaire de déterminer explicitement les conditions sur le processus \( r \) afin que les prix théoriques des obligations T \( p(\cdot, N) \) satisfassent la condition
\end{enumerate}
\begin{equation}
    p(0, N) = p^*(0, N), \quad \text{pour tout } N \leq \bar{N}.
\end{equation}

Où \( p(0, N) \)représente le prix théorique calculé dans le modèle au temps \( t=0 \) (souvent le moment présent ou le point d'évaluation), pour une obligation (ou un instrument financier similaire) qui arrive à échéance au temps \( N \).

\( p^*(0, N) \) fait référence au prix observé sur le marché au temps \( t=0 \) pour une obligation avec la même échéance \( N \), ou c'est ce qu'on appelle les données du marché.

Dans les modèles financiers, notamment dans la création de modèles de tarification, le prix théorique \( p(0, N) \) doit être aligné avec le prix du marché \( p^*(0, N) \) à travers le processus de calibrage. Ce processus assure que le modèle reflète non seulement une logique financière théorique mais aussi qu'il capte précisément les conditions actuelles du marché.

La comparaison entre \( p(0, N) \) et \( p^*(0, N) \) permet de vérifier l'efficacité et l'applicabilité du modèle. Si les différences entre ces deux sont significatives, il peut être nécessaire d'ajuster les hypothèses ou les paramètres d'entrée du modèle.

Lors de la modélisation des taux d'intérêt à court terme sur les marchés financiers, nous voulons définir la dynamique de ces taux par un processus stochastique sous la mesure de harness. Ce processus stochastique repose sur un certain nombre de paramètres qui contrôlent le comportement et les résultats du modèle.

\begin{defn}
Une mesure de martingale avec numéraire \( B \) est une mesure de probabilité \( Q \) équivalente à \( P \), par rapport à laquelle les processus de prix actualisés des obligations T sont des martingales, c'est-à-dire qu'il est vérifié que :
\[
\tilde{p}(n, N) = \mathbb{E}^Q[\tilde{p}(n + 1, N) \mid \mathcal{F}_n], \quad 0 \leq n < N \leq \bar{N}.
\]
\end{defn}

La condition de martingale est équivalente à
\begin{equation}
    p(n,N) = E^Q[D(n,N)|\mathcal{F}_n], \quad 0 \leq n < N \leq \bar{N}
\end{equation}


Le modèle court est un modèle de processus stochastique en temps discret utilisé pour attribuer une dynamique au taux d'intérêt à court terme \( r \). Dans ce modèle, la dynamique du taux est directement déterminée par une mesure martingale \( Q \). Dans la configuration du modèle, \( r \) est considéré comme un processus de Markov défini sur un espace de probabilité filtré \( (\Omega, \mathcal{F}, Q, (\mathcal{F}_n)) \), et son noyau de transition est donné par \( Q_n \).
\[
(Q_n \phi)(r_n) := \mathbb{E}^Q[\phi(r_{n+1}) \mid \mathcal{F}_n] = \int_{\mathbb{R}} \phi(\theta) Q_n (r_n, d\theta)
\]

$Q_n$ dénote deux choses différentes. A gauche de l'équation, $Q_n$ est une application linéaire qui s'applique à une fonction phi et $Q_n$ phi est également une fonction, définie par l'expression du milieu. A droite de l'équation, $Q_n$ est une mesure de probabilité, où $Q_n(r_n, A)$ = probabilité que $r_{n+1}$ appartienne à l'ensemble A, sachant $r_n$ (ou sachant $F_n$ ce qui équivalent ici). L'intégrale est donc une intégrale de Lebesgue par rapport à la mesure $Q_n$, comme lorsqu'on écrit une espérance sous la forme d'une intégrale par rapport à sa mesure de probabilité. 

pour chaque fonction intégrable à valeur réelle $\phi$ sur $\mathbb{R}$. Nous permettons à l'état de l'espace d'être $\mathbb{R}$ et donc pas nécessairement discret.

À l'échéance \( N \), la valeur nominale des obligations \( P(n, N) \) est 1. Sous une mesure risque-neutre \( Q \), la formule de valorisation de l'obligation s'exprime par :
\[
P(n, N) = \mathbb{E}^Q \left[ \exp{-\int_n^N r_u \, du} \mid \mathcal{F}_n \right].
\]
Ici, \( Q \) est une mesure de probabilité neutre au risque. En établissant et en ajustant la courbe de rendement, et en dérivant le taux forward instantané, on obtient :
\[
f(n, N) = -\frac{\partial \ln P(n, N)}{\partial N}.
\]
Le modèle d'équilibre (sans arbitrage) associe le modèle d'équilibre à l'absence d'arbitrage et la structure à terme des taux d'intérêt. Le taux d'intérêt modélisé sous forme de dérive ne permet pas à l'arbitrage de se produire. Pour \( T \), si la structure à terme des taux d'intérêt n'est pas alignée, la courbe de rendement ajustée peut ne pas être cohérente sur le long terme. Parmi les modèles les plus connus, nous trouvons le modèle de Vasicek, Cox-Ingersoll-Ross et autres.

Supposé $r$ par récursive : 
\begin{equation}
    r_{n+1} = r_n + (\Phi_n - a_n r_n) \Delta + \sigma_n \sqrt{\Delta} W_n, \quad n = 0, \ldots, \bar{N} - 1
\end{equation}
Où 
\begin{itemize}
    \item \( r_n \) représente le taux d'intérêt à court terme pour la période \( n \).
    \item \( \Phi_n \) est un paramètre dépendant du temps, utilisé pour ajuster l'objectif de la régression vers la moyenne du taux d'intérêt. Il peut être compris comme ajustant la vitesse ou la direction avec laquelle le taux d'intérêt revient à un niveau moyen.
    \item \( a_n \) est également un paramètre dépendant du temps, représentant la vitesse de régression vers la moyenne. Il détermine la rapidité avec laquelle le taux d'intérêt revient à sa valeur moyenne à long terme.
    \item \( \Delta \) est le pas de temps, c'est-à-dire l'intervalle de temps entre les périodes \( [n, n+1] \).
    \item \( \sigma_n \) représente la volatilité du taux d'intérêt pour la période \( n \), un paramètre dépendant du temps qui représente l'incertitude ou le risque de variations du taux d'intérêt.
    \item \( W_n \) est une variable aléatoire qui représente les perturbations aléatoires distribuées normalement, utilisée pour modéliser les fluctuations aléatoires du marché pendant la période \( n \).
\end{itemize}

Comme $W_n \sim \mathcal{N}(0,1)$,  par théorème central limite inverse, on a 
\begin{equation}
    r_{n+1} \sim \mathcal{N} \left( {r_n + (  \Phi_n - a_n r_n) \Delta, \sigma^2_n \Delta} \right)
\end{equation}
On introduit la forme générale de model Hull-White qui est une équation différentielle stochastique de la forme : 
$\displaystyle dX_{t}=\mu _{t}\,dt+\sigma _{t}\,dW_{t}$ qui est définie comme : 
\begin{equation*}
    dr_t = (\Phi(t)-a(t)r_t)dt +\sigma (t)dW_t
\end{equation*}
où $\sigma = \sigma (t)$ est la fonction de volatilité (déterministe), $a = a(t)$ est la vitesse (ou le taux) de retour à la moyenne, $\Phi$ est une fonction qui régit la moyenne des taux longs et $W$ est un mouvement brownien réel.

Nous insistons sur le fait que le processus r peut ici prendre des valeurs négatives arbitrairement grandes. A partir de la dynamique précédente, nous voyons maintenant comment obtenir les prix des T-bonds. A cet effet, nous utilisons la formul \(2.2\) qui donne l'expression des prix p(n,N) en termes d'espérance conditionnelle sous Q des taux courts $r_k$ avec $n \le k < N$. Avant d'énoncer le résultat suivant, nous introduisons quelques notations que nous utiliserons systématiquement dans la suite. Mettre


Nous commençons définir le prix des T-bonds une fonction $\varphi$, soit 
\begin{equation}
    \varphi_0 (r) = e^{-r}
\end{equation}

et $\varphi_{n}^{N}$ tel que 

\begin{equation}
\begin{cases}
    \varphi^{N}_{N-1}(r) = 1\\
\varphi^{N}_{n-1}(r) = Q_{n-1}(\Phi^{N}_n) (r), \quad 1 \leq n \leq N\\
\end{cases}
\end{equation}

pour tout $N \le \bar{N}$ et $r \in \mathbb{R}$

Soit la mesure martingale $Q$,
\begin{equation}
p(n, N) = e^{-r_n} \phi_n^N (r_n), \quad 0 \leq n < N \leq \overline{N},
\end{equation}
Car 
\begin{equation*}
    \frac{p(n,N)}{p(n,n+1)} = E^Q[p(n+1,N)|F_n], \quad 0 \leq n < N \leq \bar{N}
\end{equation*}
on a 
\begin{align*}
p(n-1,N) 
    &= p(n-1,n) E^Q[p(n,N)|\mathcal{F}_{n-1}]\\
    &= e^{-r_{n - 1}} \mathbb{E}^Q \left[ e^{-r_n} \phi_n^N (r_n) \mid \mathcal{F}_{n-1} \right] \\
    &= e^{-r_{n - 1}} Q_{n-1} (\phi_n^N (r_{n-1})) = e^{-r_n - 1} \phi_{n-1}^N (r_{n-1})\\
\end{align*}


\subsection{Modèle Affine}

Si des fonctions \( f_n, g_n \) existent telles que
\begin{equation}
\tilde{Q}_n(r, \lambda) = \exp(-f_n(\lambda) - g_n(\lambda)r), \quad r, \lambda \in \mathbb{R},
\label{eq:my_equation_2.8}
\end{equation}
pour chaque \( n, 0 \leq n \leq N \), alors les fonctions \( \varphi_n^N \) dans (4.18) ont l'expression suivante :
\begin{equation}
\varphi_n^N (r) = \exp(-A_n^N - B_n^N r), \quad r \in \mathbb{R},
\label{eq:my_equation_2.9}
\end{equation}
avec les constantes \( A_n^N, B_n^N \) définies par les récursions suivantes :
\begin{equation}
\begin{cases}
    A_{N-1}^N=B_{N-1}^N=0, \\
    A_{n-1}^N=A_n^N+f_{n-1}\left(1+B_n^N\right), \\
    B_{n-1}^N=g_{n-1}\left(1+B_n^N\right) 
\end{cases}    
\end{equation}

\begin{defn}
    Nous disons qu'un modèle court, les prix des T-bonds $p(n,N)$ est un modèle de structure par termes affine.
    \begin{equation}
    p(n, N) = \exp(-A_n^N - (1 + B_n^N) r_n), \quad 0 \leq n < N \leq \overline{N}.
    \label{eq:my_equation_2.11}
    \end{equation}
    
\end{defn}

\subsection{Modèle de Hull-White à temps discret}

Rappelons que 
\begin{equation}
    r_{n+1} = r_n + (\Phi_n - a_n r_n) \Delta + \sigma_n \sqrt{\Delta} W_n, \quad n = 0, \ldots, \bar{N} - 1
\end{equation}
Pour la simplicité, nous ne considérons que $a_n$ et $\sigma_n$ sont constantes, à savoir $a_n \equiv a $ et $ \sigma_n \equiv \sigma$, donc on a 
\begin{equation}
    r_{n+1} = r_n + (\Phi - a r_n) \Delta + \sigma \sqrt{\Delta} W_n, \quad n = 0, \ldots, \bar{N} - 1
\end{equation}

Soit la fonction génératrice des moments pour une variable aléatoire normalement distribuée \( X \sim \mathcal{N}(\mu, \sigma^2) \) est donnée par :
\begin{equation}
M_X(t) = \exp\left(\mu t + \frac{\sigma^2 t^2}{2}\right)
\end{equation}

Alors
\begin{equation}
\tilde{Q}_n(r_n, \lambda) = \exp(-f_n(\lambda) - g_n(\lambda)r), \quad r, \lambda \in \mathbb{R},
\end{equation}
avec 
\begin{equation}
    f_n{\lambda} = \Phi_n \Delta \lambda -\frac{\sigma^2\Delta}{2}\lambda^2 \quad et \quad g_n (\lambda) = \lambda(1-a\Delta)
\end{equation}
donc on a par \ref{eq:my_equation_2.8},\ref{eq:my_equation_2.9},\ref{eq:my_equation_2.11}, on a 
\begin{equation}
    p(n, N) = \exp(-A_n^N - (1 + B_n^N) r_n), \quad 0 \leq n < N \leq \overline{N}.
    \end{equation}
où
\begin{equation}
    \begin{cases}
    A_{N-1}^{N} = B_{N-1}^{N} = 0, \\
    A_{n-1}^{N} = A_n^{N} + \Phi_{n-1} \Delta (1 + B_n^{N}) - \frac{\sigma^2 \Delta}{2} (1 + B_n^{N})^2, \\
    B_{n-1}^{N} = (1 + B_n^{N})(1 - a\Delta).
    \end{cases}
\end{equation}
pour tout $n \le N - 2 $

\begin{equation}
    \begin{cases}
        A_n^N = \sum_{k=n}^{N-2} \left( \Phi_k \Delta (1 + B_{k+1}^N) - \frac{\sigma^2 \Delta}{2} (1 + B_{k+1}^N)^2 \right), \\
    B_n^N = (1 - a\Delta)^{N-n-2} (N - n - 1 - a\Delta).
    \end{cases}
\end{equation}

dans le cas simplifié $a = 0$, la relation est donnée par

\begin{align}
A_n^N - A_{n+1}^N &= (N - n - 1)\Phi_n \Delta - \frac{(N - n - 1)^2 \sigma^2 \Delta}{2}, \\
(1 + B_n^N)r_n - (1 + B_{n+1}^N)r_{n+1} &= (N - n)r_n - (N - n - 1)(r_n + \Phi_n \Delta + \sigma \Delta W_n),
\end{align}
et le taux \( r_n \) 
\begin{equation}
r_n = -\frac{\log p(n, N) + A_n^N}{N - n},
\end{equation}
par la recursivité, le prix des T-bonds donné pars:
\begin{equation}
p(n + 1, N) = p(n, N) \exp\left( -\frac{\log p(n, N) + A_n^N}{N - n} - \frac{\sigma^2 \Delta (N - n - 1)^2 - (N - n - 1)\sigma \Delta W_n}{2} \right).
\end{equation}

Dans le modèle de Hull-White en temps discret,
\[
r_{n+1} = r_n + \Phi_n \Delta + \sigma \Delta W_n,
\]
on a 
\[
\Phi_n = \frac{R(0,n+1) - R(0,n)}{\Delta} + \sigma^2 \left(n + \frac{1}{2}\right).
\]

 Pour chaque choix de \(\sigma\), la calibration du modèle consiste simplement à poser
\[
\Phi_n = \frac{R^*(0,n+1) - R^*(0,n)}{\Delta} + \sigma^2 \left(n + \frac{1}{2}\right),
\]
où \(R^*\) est le taux forward du marché, où
\[
R^*(0,n) = \log \left(\frac{p^*(0,n)}{p^*(0,n+1)}\right).
\]


%----------------------------------------------------------------------------------------
%	SECTION
%----------------------------------------------------------------------------------------
Nous décrivons ensuite les modèles de type "forward", pour lesquels certains des problèmes mentionnés ci-dessus ne se posent pas.
\section{Modèles Forward}

Dans un modèle forward, on attribue directement la dynamique des processus de prix $p(·,N)$ pour chaque $N=1,...,\bar{N}$. Dans ce cas, la structure à terme initiale $p^*(0,N), N = 1,...,\bar{N}$, est automatiquement supposée comme la donnée initiale des processus de prix.
Le résultat suivant contient une caractérisation de la mesure de martingale exprimée en termes d'une relation importante entre les prix des $T$-bonds :

\setcounter{equation}{0}
\renewcommand{\theequation}{\thesection.\arabic{equation}}
\begin{equation}
    \frac{p(n,N)}{p(n,n+1)} = E^Q[p(n+1,N)|F_n], \quad 0 \leq n < N \leq \bar{N}
    \label{eq:my_equation_2.2.1}
\end{equation}


En examinant la formule (2.2.1), on introduit une dynamique récursive du type suivant

\begin{equation}
    p(n,N) = \frac{p(n-1,N)}{p(n-1,n)} \mu_{n,N}, \quad 0 \leq n < N \leq \bar{N}
    \label{eq:my_equation_2.2.2}
\end{equation}


Où $\mu_{n,N}$ sont des variables aléatoires définies sur un espace de probabilité filtré $(\Omega,F,\mathbb{P},(\mathcal{F}_n))$. D'après (\ref{eq:my_equation_2.2.1})et (\ref{eq:my_equation_2.2.2}) on obtient la formule suivante, d'où implique que les prix actualisés des $T$-bonds sont des martingales dans la mesure de probabilité $Q$, elle constitue la base de l'étude de l'existence et de l'unicité de la mesure de martingale (ou des propriétés d'absence d'opportunité d'arbitrage et de complétude du modèle).

\begin{equation}
    E^Q[\mu_{n,N}|F_{N-1}] = 1, \quad 0 \leq n < N \leq \bar{N}
    \label{eq:my_equation_2.2.3}
\end{equation}

Preuve.

D'après (2.2.2), $\mu_{n,N}={p(n,N)p(n-1,n) \over p(n-1,N)}$, alors on a :

\begin{align*}
    E[\mu_{n,N} | F_{n-1}] 
    & = E\left[\frac{p(n,N)p(n-1,n)}{p(n-1,N)} | F_{n-1}\right] \\
    & = E[p(n,N) | F_{n-1}] \frac{p(n-1,n)}{p(n-1,N)} \\
    & = \frac{p(n-1,N)}{p(n-1,n)} \frac{p(n-1,n)}{p(n-1,N)} \text{(car} \ {p(n-1,n) \over p(n-1,N)} \ \text{est} \ F_{n-1}\ \text{mesurable)}\\
    & = 1 \tag*{$\Box$}
\end{align*}



Le résultat suivant donne l'expression explicite des processus de prix en termes de données initiales
\begin{equation}
    p(n,N) = \frac{p(0,N)}{p(0,n)} \prod_{k=1}^{n}\frac{\mu_{k,N}}{\mu_{k,n}}
    \label{eq:my_equation_2.2.4}
\end{equation}


Remarquons que (\ref{eq:my_equation_2.2.4}) exprime le prix d'un T-bonds en termes des prix initiaux $p(0,\cdot)$ ( observable sur le marché ) et des facteurs stochastiques $\mu_{n,N}$ du modèle considéré. Par conséquent, comme nous l'avons déjà souligné, une caractéristique commode des modèles forward est le fait que la structure initiale des termes $p^*(0,N)$, $N = 1, . . . , \bar{N}$ , est automatiquement supposée comme la donnée initiale du processus de prix.

En combinant la formule récursive (2.2.2) pour les prix des $T$-bonds avec la définition (1.5) des taux, nous obtenons la formules récursive suivante pour les processus des taux.

$$L(n,N)=(L(n-1,N)+{1 \over \triangle}) {\mu_{n,N} \over \mu_{n,N+1}} - {1 \over \triangle}$$
    \text{Preuve.}

\begin{align*}
    L(n,N) &= \frac{1}{\triangle} \left( \frac{p(n,N)}{p(n,N+1)} - 1 \right) (ici, M-N=N+1-N=1)\\
    &= \frac{1}{\triangle} \left( \frac{p(n-1,N)\mu_{n,N}}{p(n-1,n)p(n,N+1)} - 1 \right) (\text{d'après} (2.2.2))\\
    &= \frac{1}{\triangle} \left( \frac{p(n-1,N)\mu_{n,N} p(n-1,n)}{p(n-1,n)p(n-1,N+1) \mu_{n,N+1}} - 1 \right) (\text{d'après} (2.2.2))\\
    &= \frac{1}{\triangle} \left( \frac{p(n-1,N)\mu_{n,N}}{p(n-1,N+1)\mu_{n,N+1}} - 1 \right) \\
    &= \left( L(n-1,N) + \frac{1}{\triangle} \right) \frac{\mu_{n,N}}{\mu_{n,N+1}} - \frac{1}{\triangle} \tag*{$\Box$}
\end{align*}



De la même manière, on obtient les formules suivantes : 

$$R(n,N)=R(n-1,N) + \log({\mu_{n,N} \over \mu_{n,N+1}})$$

$$r_n=-r_{n-1}-\log(p(n-1,n+1))-\log(\mu_{n,n+1})$$

Le modèle binomial forward est une extension du modèle binomial classique utilisé en finance pour évaluer les options et autres instruments dérivés. Dans ce modèle, on considère la dynamique des prix d'un actif financier sur une période de temps discrète, en supposant que le prix de l'actif peut soit augmenter, soit diminuer à chaque étape, avec des probabilités fixées.

Dans le contexte des taux d'intérêt et des obligations, le modèle binomial forward peut être utilisé pour estimer l'évolution future des taux d'intérêt et des prix des obligations à différentes échéances. Plus précisément, on peut utiliser ce modèle pour estimer le prix d'une obligation ou d'un produit dérivé lié aux taux d'intérêt dans le futur, en supposant des mouvements binomiaux des taux d'intérêt.

Ici, on considère un modèle forward du type (2.2.1) dans lequel \(\mu_{n,N}\) sont des variables aléatoires qui peuvent prendre deux valeurs :
\begin{equation}
\mu_{n,N} =
\begin{cases}
u_{n,N} \\
d_{n,N}
\end{cases} 
\end{equation}

où \(0 < d_{n,N} < u_{n,N}\). Si \(d_{n,N} < 1 < u_{n,N}\), alors il existe une mesure de martingale \(Q\). \(Q\) est une mesure de martingale si et seulement si (2.2.3) tient en compte, alors on a :
$$
1=E^Q[\mu_{n,N}|F_{N-1}]=u_{n,N}Q(\mu_{n,N}=u_{n,N}|F_{N-1})+d_{n,N}Q(\mu_{n,N}=d_{n,N}|F_{N-1})
$$
d'où
    \begin{align*}    
    q_{n,N} :&= Q(\mu_{n,N} = u_{n,N})\\
    &= 1 - Q(\mu_{n,N} = d_{n,N}) \\
    &= \frac{1 - d_{n,N}}{u_{n,N} - d_{n,N}} 
    \end{align*}

Nous prenons un exemple pour bien comprendre. On considère un modèle binomial forward avec \(\bar{N} = 3\) et suppose la dynamique (2.2.1) avec \(u_{n,N} \equiv u = 2\) et \(d_{n,N} \equiv d = \frac{1}{2}\) pour \(1 \leq n \leq N \leq 3\).

D'abord on exprime les prix des $T$-bonds dans le modèle binomial forward avec $\bar N=3$ : 
\begin{align*}
\text{Si } N &= 1, \text{ alors le prix de } T\text{-bonds est } p(0,1) \text{ à } t=0 \text{ et } 1 \text{ à } t=1; \\
&\text{Ici}, \text{p(0,1) est considéré comme donné initial de T-bonds}\\
\text{Si } N &= 2, \text{ alors le prix de } T\text{-bonds est } p(0,2) \text{ à } t=0, \text{ } p(1,2) = \frac{p(0,2)}{p(0,1)}\mu_{1,2} \text{ à } t=1, \text{ et } 1 \text{ à } t=2; \\
&\text{d'après (2.2.2) à t=1, et 1 à t=2}\\
\text{Si } N &= 3, \text{ alors le prix de } T\text{-bonds est } p(0,3) \text{ à } t=0, \text{ } p(1,3) = \frac{p(0,3)\mu_{1,3}}{p(0,1)\mu_{1,1}} \text{ à } t=1, \\
& \quad p(2,3) = \frac{p(0,3)}{p(0,2)}\frac{\mu_{1,3}\mu_{2,3}}{\mu_{1,2}} \text{ à } t=2, \text{ et } 1 \text{ à } t=3;
\end{align*}
On fait le calcul de $p(2,3)$ comme suit :
\begin{align*}
    p(2,3) 
    &= {p(1,3) \over p(1,2)}\ \mu_{2,3}\\
    &={{p(0,3) \over p(0,1)}\ \mu_{1,3} \over {p(0,2) \over p(0,1)}\mu_{1,2}}\ \mu_{2,3}  \quad \text{  (d'après\ (2.2.2))}\\
    &={p(0,3) \over p(0,2)}{\mu_{1,3}\mu_{2,3} \over \mu_{1,2}}
\end{align*}

alors \(w = (\mu_{1,2}, \mu_{1,3}, \mu_{2,3})\), avec \(\mu_{1,2}, \mu_{1,3}, \mu_{2,3} \in \{u, d\}\). Donc l'espace de probabilité contient 8 éléments :
\begin{align*}
w_1 &= (u, u, u), & w_2 &= (u, u, d), & w_3 &= (u, d, u), & w_4 &= (u, d, d), \\
w_5 &= (d, u, u), & w_6 &= (d, u, d), & w_7 &= (d, d, u), & w_8 &= (d, d, d).
\end{align*}

Ensuite, on note \(q_k = Q(\{w_k\})\) pour \(k = 1, \ldots, 8\). Alors on obtient :
\[
Q(\mu_{1,2} = u) = \frac{1 - d_{1,2}}{u_{1,2} - d_{1,2}} = \frac{1 - \frac{1}{2}}{2 - \frac{1}{2}} = \frac{1}{3} \quad \text{d'après\ (2.2.6)}
\]

ensuite, \(Q(\mu_{1,2} = u) = Q(\mu_{1,3} = u) = Q(\mu_{2,3} = u)\), \(\mu_{1,2}, \mu_{1,3}, \mu_{2,3}\) sont deux à deux indépendantes, par 
\begin{align*}
    Q(\mu_{1,2}=u)&=q_1+q_2+q_3+q_4=\frac{1}{3}, \\
    Q(\mu_{1,3}=u)&=q_1+q_2+q_5+q_6=\frac{1}{3}, \\
    Q(\mu_{2,3}=u)&=q_1+q_3+q_5+q_7=\frac{1}{3}
\end{align*},on obtient le système suivant :
\[
\begin{cases}
q_1 + q_2 + q_3 + q_4 = \frac{1}{3}, \\
q_1 + q_2 + q_5 + q_6 = \frac{1}{3}, \\
q_1 + q_3 + q_5 + q_7 = \frac{1}{3}.
\end{cases}
\]
De plus, \(\mu_{2,3}\) est indépendante de \(\mu_{1,3}\) et \(\mu_{1,2}\), on a 
\begin{align*}
q_1&=Q(\mu_{1,2}=\mu_{1,3}=u)\ Q(\mu_{2,3}=u)=\frac{1}{3}(q_1+q_2)\\
q_4&=Q(\mu_{1,2}=u,\mu_{1,3}=d)\ Q(\mu_{2,3}=d)=\frac{2}{3}(q_3+q_4)\\
q_5&=Q(\mu_{1,2}=d,\mu_{1,3}=u)\ Q(\mu_{2,3}=u)=\frac{1}{3}(q_5+q_6)\\
q_7&=Q(\mu_{1,2}=\mu_{1,3}=d)\ Q(\mu_{2,3}=u)=\frac{1}{3}(q_7+q_8)\\    
\end{align*}

d'où \(q_2 = 2q_1\). De même, \(q_4 = 2q_3, q_6 = 2q_5, q_8 = 2q_7\). Avec le système précédant, on obtient :
\[
\begin{cases}
q_1 = q - \frac{1}{9}, \\
q_2 = 2q_1, \\
q_3 = \frac{2}{9} - q = q_5, \\
q_4 = 2q_3, \\
q_6 = 2q_5, \\
q_7 = q, \\
q_8 = 2q_7.
\end{cases}
\]
avec \(\frac{1}{9} < q < \frac{2}{9}\). En prenant par exemple \(q = \frac{4}{27}\), nous obtenons l'unique mesure martingale par rapport à laquelle \(\mu_{1,2}, \mu_{1,3}, \mu_{2,3}\) sont des variables aléatoires indépendantes.




%\input{chapters/article}
%\input{chapters/chaptercontrib}

%\part*{Discussion et conclusion générale}
% \fancyhead{} % clear all header fields
% \fancyhead[OL]{\textsc{Discussion et perspectives}}
% \input{chapters/discussion.tex}


\chapter{Conclusion} % Main chapter title
%\addcontentsline{toc}{chapter}{Conclusion générale}  
%\startcontents[chapters]
%\pagestyle{plain} % remove headers/footers from the chapt\markboth{Left}{Right}
 


%\pagestyle{fancy}
%\renewcommand{\chaptermark}[1]{\markboth{Chapter \thechapter. #1}{}}
%\renewcommand{\sectionmark}[1]{\markright{\thesection\ #1}}


Dans ce rapport, nous avons exploré en profondeur le marché des taux d'intérêt, un domaine crucial de la finance qui influence de manière significative de nombreux aspects de l'économie mondiale. En commençant par une définition claire du marché des taux d'intérêt,  nous avons établi les bases nécessaires à la compréhension de ce domaine complexe.

En classant les taux d'intérêt en différentes catégories sur la base de leurs caractéristiques, nous avons offert un aperçu approfondi des différentes dimensions de ce marché dynamique. Cette analyse nous a permis de mieux saisir l'ampleur de l'influence des taux d'intérêt sur les décisions d'investissement et de prêt.

En détaillant les modèles stochastiques de type "short" et "forward", notamment les modèles affine, Hull-White et binomial forward, nous avons fourni des outils puissants pour la modélisation et l'analyse des taux d'intérêt. Ces modèles nous permettent d'identifier et de comprendre les principales forces qui façonnent l'évolution des taux d'intérêt, ainsi que les implications économiques qui en découlent.

En conclusion, le marché des taux d'intérêt est un domaine complexe mais essentiel de la finance, qui influence de manière significative l'ensemble de l'économie mondiale. En comprenant les mécanismes et les tendances qui le sous-tendent, les acteurs du marché peuvent prendre des décisions éclairées et contribuer à la stabilité et à la croissance économique à long terme.


%----------------------------------------------------------------------------------------
%	BIBLIOGRAPHY
%----------------------------------------------------------------------------------------
%\printbibliography %Prints bibliography
\fancyhead{} % clear all header fields
\fancyhead[OL]{\textsc{Bibliographie}}
\printbibliography[
heading=bibintoc,
title={Bibliographie}]

%----------------------------------------------------------------------------------------

%: ----------------------- glossary ------------------------
%\printindex
%\fancyhead{} % clear all header fields
%\fancyhead[OL]{\textsc{Glossaire}}
%\printglossary[title=Glossaire,toctitle=Glossaire]

%\input{auxilliaires/glossaire} 


\end{document}  
